% 引言章节
% 说明：
% 1. 本文件设计为被主文档 \input{} 进来使用。
% 2. 内容来源于 04_paper.md 的 Introduction 部分。
% 3. 这里只做 Markdown 到 LaTeX 的语法迁移，不改写正文内容。

\section{Introduction}

增量式三维场景重建在灾害救援、自动驾驶与工业巡检等时效性强的现场任务中具有重要需求。近年来，LiDAR-惯导-视觉融合的增量式高斯建图系统已在大尺度室外场景中展现出在线照片级重建的能力：外部前端提供稳定位姿与多模态观测，后端以高斯原语逐帧扩展地图。然而，在线系统每帧的优化预算有限，新插入的高斯往往仅经历少量优化即需服务于后续帧的渲染，其初始属性质量直接影响地图的累积建图效果。

现有系统在地图后端的表示与初始化两方面仍存在相互关联的不足。表示层面，标准 3DGS 以三维椭球为基础原语，缺乏显式的表面法线语义，限制了几何先验与地图表示的直接耦合；2DGS 已通过将高斯压缩为二维有向盘面在几何质量上取得显著进步，但尚未被引入增量式建图。初始化层面，现有系统普遍采用基于规则的策略为新高斯赋予属性，导致新高斯需要较多优化步数收敛，与有限的逐帧预算矛盾；虽然已有系统引入深度补全缓解 LiDAR 覆盖不足，但学习型先验仅用于确定补点位置，尚未扩展到属性全面预测。若表示缺乏几何可解释性，规则初始化便难以有效利用外部先验；而缺少学习型属性预测，表面化表示的优势也无法在插入阶段充分兑现。

针对上述问题，本文提出一种面向增量式 2DGS 场景重建的上下文感知前馈高斯回归系统。系统以 2DGS surfel 为地图表示，在外部 LIVO 前端提供位姿的条件下，融合深度补全与单目法线先验构造候选点及其几何方向；前馈回归网络聚合当前帧与滑动窗口内历史帧的图像-渲染图上下文，在单次前向传播中直接预测新高斯的完整 surfel 属性，随后纳入在线优化。网络训练通过我们构建的一套自蒸馏训练方法完成，无需外部标注数据集。

本文的主要贡献如下：

\begin{itemize}
    \item 基于 2DGS surfel 的几何感知增量式高斯地图后端。不同于以 3DGS 体高斯为默认表示的现有增量系统，本文引入表面化的 2DGS 表示，使地图天然具备显式法线语义，从而能够直接耦合深度补全与法线先验，为前馈初始化提供几何可解释的表示基础。

    \item 上下文感知的前馈高斯插入方法。不同于依赖规则逐一赋值的传统初始化方式，本文通过前馈网络以逐点跨注意力融合多视角上下文,并以法线先验为锚点，在单次前向传播中直接预测完整 surfel 属性，使新高斯在插入时就更接近收敛状态，显著减少有限优化预算下的早期质量损失。

    \item 面向在线增量建图的自蒸馏训练。前馈网络无需外部标注数据集，仅从系统自身的增量建图过程中采集监督信号，实现从在线建图到离线训练再到重新部署的自增强过程。
\end{itemize}
